\documentclass[11pt,a4paper]{article}

% Compile with pdfLaTeX (the Overleaf default).
\usepackage[utf8]{inputenc}
\usepackage[T1]{fontenc}
\usepackage[english]{babel}
\usepackage{lmodern}
\usepackage{microtype}
\usepackage[a4paper,margin=2.3cm]{geometry}
\usepackage{amsmath}
\usepackage{graphicx}
\usepackage{booktabs}
\usepackage{tabularx}
\usepackage{array}
\usepackage{float}
\usepackage{caption}
\usepackage{xcolor}
\usepackage{listings}
\usepackage[hidelinks]{hyperref}

\graphicspath{{figures/}}
\captionsetup{font=small,labelfont=bf}
\setlength{\parskip}{5pt}
\setlength{\parindent}{0pt}
\renewcommand{\arraystretch}{1.15}

\definecolor{wrong}{HTML}{B42828}
\definecolor{notebg}{HTML}{F3F2EE}
\newcommand{\WRONG}{\textcolor{wrong}{\textbf{WRONG}}}
\newcommand{\OK}{\textbf{OK}}
\newcommand{\timesx}{\ensuremath{\times}}
% Red marker for every value that still has to be filled in from your own results.
% Search the file for TODO before handing the report in.
\newcommand{\TODO}[1]{\textcolor{wrong}{\textbf{[TODO: #1]}}}
% Include a figure if the file exists, otherwise show a visible placeholder.
\newcommand{\fig}[2]{\IfFileExists{figures/#1}{\includegraphics[width=#2]{#1}}%
  {\fbox{\parbox{0.9\linewidth}{\centering\vspace{2em}\TODO{figure \detokenize{#1} not generated yet}\vspace{2em}}}}}
\newcommand{\question}[1]{\par\medskip\textbf{#1}\par}
\newcommand{\notebox}[1]{\par\smallskip\noindent\colorbox{notebg}{\parbox{\dimexpr\linewidth-2\fboxsep}{\small #1}}\par\smallskip}

\lstset{basicstyle=\ttfamily\small,columns=fullflexible,keepspaces=true,
        frame=single,rulecolor=\color{black!25},backgroundcolor=\color{notebg},
        xleftmargin=3pt,xrightmargin=3pt,breaklines=true}

\begin{document}

% Title
\begin{center}
  {\small NETWORK PROGRAMMING \textperiodcentered{} UTM}\\[6pt]
  {\LARGE\bfseries Laboratory Work 1: Concurrency in Practice}\\[10pt]
  \begin{tabular}{rl}
    Authors: & Pancenco Ina, Rusnac Nichita \\
    Group:   & FAF-242 \\
    Teacher: & Ernest B\^{\i}tca \\
    Date:    & October 2026
  \end{tabular}
\end{center}
\vspace{4pt}\hrule\vspace{8pt}

% Section 1
\section{Hardware and software}

All measurements in this report were taken on a single laptop. Its configuration is listed in the
table below, because the number and kind of processor cores, and the speed of the storage devices,
directly determine how the different tasks behave.

\begin{center}
\small
\begin{tabularx}{\linewidth}{@{}lX@{}}
\toprule
Component & Value \\
\midrule
CPU & AMD Ryzen 7 7435HS (8 identical cores, each with simultaneous multithreading) \\
Physical cores / logical processors & 8 / 16 \\
Clock & 3.1\,GHz base (as reported by \texttt{lab1 info}), up to 4.5\,GHz boost \\
RAM & 16\,GB DDR5, one Samsung module (M425R2GA3PB0-CWMOD), running at 4\,800\,MT/s \\
System disk & NVMe SSD, Micron MTFDKBA1T0QFM, 1\,TB \\
USB flash drive (Task~2) & Transcend JetFlash, 4\,GB, FAT32; measured sequential read about 32\,MB/s
  with an empty cache, which is the speed class of a USB~2.0 drive \\
USB port & AMD USB 3.10 eXtensible host controller, USB 3.0 root hub \\
Operating system & Windows 11 Pro, version 26H2 (build 26300) \\
Language and runtime & C\# on .NET 9 (runtime 9.0.20, x64), Release build \\
Power & AC power, power plan ``Turbo'' \\
\bottomrule
\end{tabularx}
\end{center}

The processor shapes almost every result of Task~1. The Ryzen 7 7435HS has 8 physical cores, and every
core can run two hardware threads at the same time (simultaneous multithreading, SMT, called
Hyper-Threading by Intel). The operating system therefore sees 8 physical cores but 16 logical
processors. The two hardware threads of one core are not two full cores: they share the core's
execution units and caches, and the second thread mainly uses the moments in which the first one
waits, for example for memory. As lecture~1 explained, the second thread of a core therefore adds
only a modest gain, in the order of 15 to 30 per cent, and nothing like a second core. All 8 cores
are of the same kind and run at the same speed, so equal parts of work take about the same time on
every core. The two numbers 8 and 16 are the reference lines in the charts of Task~1.

% The note about the faulty "MyMedia" flash drive belonged to the other laptop's measurements.
% Re-enable it (and adapt it) only if your own Task 2 drive had the same problem.
%\notebox{\textbf{A note on the USB drive.} ...}

% Section 2
\section{Method}

\textbf{The program.} The whole laboratory work is implemented as a single C\# console program. Every
task and every configuration is selected by command-line parameters in the form
\texttt{Labs.exe lab1 <command> [options]}; for example, Task~1a with eight threads is started with
the command \texttt{lab1 1a} and the option for eight threads. The full list of commands is given in
Appendix~\ref{app:run}.

\textbf{The input data.} The command \texttt{lab1 gen} writes the file \texttt{numbers.txt}, which
contains 50\,000\,000 random integers between $-1\,000\,000$ and $1\,000\,000$, one number per line.
The file is 352\,MB in size. The random generator is started with a fixed seed (12345), so the same
file can be produced again at any time. While the generator writes the file, it also counts how many
of the numbers are negative, zero and positive, and stores these reference counts in a second file,
\texttt{numbers.ref.txt}:

\begin{center}
\small
\begin{tabular}{rrrr}
\toprule
negative & zero & positive & total \\
\midrule
25\,007\,776 & 22 & 24\,992\,202 & 50\,000\,000 \\
\bottomrule
\end{tabular}
\end{center}

For Task~2, the command \texttt{lab1 split} cuts \texttt{numbers.txt}, keeping the original order,
into 8 files of 6\,250\,000 lines each, which are about 44\,MB per file.

\textbf{The correctness check.} After every single run, the program compares the three counts it has
computed with the reference counts and prints \OK{} if they are identical or \WRONG{} if they are
not. A configuration only counts as a valid result if it reports \OK. Every run in this report
passed this check, with the single, intended exception of the unsynchronised version of Task~1b,
whose purpose is precisely to show a wrong result.

\textbf{Threads.} All threads used by the tasks are dedicated operating-system threads that the
program creates itself with \texttt{new Thread}, starts with \texttt{Start} and waits for with
\texttt{Join}. The .NET thread pool is not used, so the number of threads in each experiment is
exactly the number that was requested.

\textbf{What is measured.} The execution time is the wall-clock time of the actual work, measured
with a \texttt{Stopwatch}. It does not include the start-up of the process, and in Task~1b it also
does not include the loading of the numbers into memory, which the assignment explicitly excludes.
The memory is reported as two high-water marks of the process: the peak working set, which is the
largest amount of physical RAM the process occupied at any moment, and the peak private memory,
which is the largest amount of memory the process had committed for itself. Because these values
are maxima over the whole lifetime of a process, every configuration was run in a fresh process, so
that one measurement cannot inherit the memory peak of another. Part of the working set belongs to
the .NET runtime itself and is present even before any work starts.

\textbf{How Task~1 was run.} Before the measured runs, one unrecorded warm-up run read the whole of
\texttt{numbers.txt}, so that the file was already held in the operating system's file cache. All
measured runs of Task~1 therefore read the file from RAM, and they measure the CPU work of parsing
and counting rather than the speed of the SSD. Every configuration was run three times, and the
report uses the median of the three times. The repetitions formed the outer loop of the measuring
script, so that a temporary slowdown of the laptop, for example a background task or the processor
getting hot, would affect one repetition of every configuration rather than all three repetitions of
a single one.

\textbf{How Task~2 was run.} Before every run of Task~2, the file cache of the USB drive was emptied by
dismounting its volume with the script \texttt{dismount-volume.ps1}, which locks the volume and sends
it the \texttt{FSCTL\_DISMOUNT\_VOLUME} request. When a volume is dismounted, Windows discards all of its
cached data and mounts it again on the next access, so this has the same effect as ejecting,
unplugging and replugging the drive, and all data really had to be read from the flash memory. A
quick test confirmed it: reading one 44\,MB part took 0.03\,s while it was cached and 1.47\,s right
after a dismount. Each configuration of Task~2 was run three times, more than the single run the
assignment requires, and the report uses the median.

\textbf{Runtime settings.} The program uses the workstation garbage collector in non-concurrent mode,
so that the runtime does not start additional garbage-collection threads that would compete with
the measured threads for the processor. Tiered compilation is switched off, so that all methods are
compiled with full optimisation from the start and short runs are not partly measured in slower,
unoptimised code.

% Section 3
\section{Task 1: One file}

\subsection{Task 1a: each thread reads and parses its own part}

In Task~1a, the file is divided into as many equal byte ranges as there are threads. Each thread
opens the file on its own, reads only its own range through a private 1\,MB buffer, and parses the
digits directly from the bytes, without creating a string for each line. It counts the negative,
zero and positive numbers in its own local counters, and the main thread adds the counts of all
threads together after joining them.

The difficult part is the boundary between two neighbouring ranges, because a boundary computed from
the file size will usually fall in the middle of a line. To solve this, both the start and the end
of every range are moved forward to the beginning of the next line. Since two neighbouring parts
apply the same rule to the same offset, the end of one part is always exactly the start of the next
one. As a result every line belongs to exactly one part, so no line is counted twice and none is
skipped.

\begin{figure}[H]
  \centering
  \fig{task1a.png}{\linewidth}
  \caption{Task 1a: execution time, speed-up and peak private memory against the number of threads.
           The dashed vertical lines mark the 8 physical cores and the 16 logical processors.}
\end{figure}

\begin{center}
\small
\begin{tabular}{lrrrrrrr}
\toprule
Threads & 1 & 2 & 4 & 8 & 16 & 32 & 64 \\
\midrule
Time, ms (median)          & 949 & 531 & 295 & 207 & \textbf{150} & 145 & 158 \\
Speed-up                   & 1.00 & 1.79 & 3.21 & 4.58 & \textbf{6.32} & 6.55 & 6.01 \\
Efficiency (speed-up / threads) & 100\% & 89\% & 80\% & 57\% & 40\% & 20\% & 9\% \\
Peak private, MB           & 8 & 9 & 11 & 15 & 24 & 42 & 73 \\
Peak working set, MB       & 26 & 27 & 29 & 33 & 42 & 58 & 90 \\
Check                      & OK & OK & OK & OK & OK & OK & OK \\
\bottomrule
\end{tabular}
\end{center}

\subsection{Task 1b: three ways to count}

In Task~1b, all 50\,000\,000 numbers are first loaded into an integer array, and this step is not
measured. The array is then divided into 16 equal parts, one for each of the 16 logical processors,
and 16 threads count the signs of their part. The three versions required by the assignment differ
only in where the counters live and how they are protected. In the unsynchronised version, all
threads increment the same three shared counters with a plain \texttt{counter++}. In the locked
version, they use the same shared counters, but every single increment happens inside a
\texttt{lock} statement. In the local version, each thread counts in its own local variables, and
the 16 partial results are added together only after all threads have finished. As an additional
fourth version, not required by the assignment, the shared counters were also incremented with the
atomic operation \texttt{Interlocked.Increment}.

\begin{center}
\small
\begin{tabular}{lrcr}
\toprule
Version & Time (median of 3) & Totals & Peak memory \\
\midrule
local (per-thread counters)        & \textbf{77\,ms} & correct & $\approx$218\,MB \\
unsync (shared, no synchronisation) & 302\,ms & \WRONG & $\approx$218\,MB \\
interlocked (shared, atomic, extra) & 1\,057\,ms & correct & $\approx$218\,MB \\
lock (shared, every increment locked) & 10\,261\,ms & correct & $\approx$218\,MB \\
\bottomrule
\end{tabular}
\end{center}

As the assignment requires, the unsynchronised version was run three times. Its totals are wrong in
every run, and they are different every time. Across the three runs, only between 7.2 and 9.5 million
of the 50 million numbers were counted, which means that between about 81 and 86 per cent of all
increments were lost. The individual runs are:

\begin{center}
\small
\begin{tabular}{crrrrrr}
\toprule
Run & Time & negative & zero & positive & total & lost increments \\
\midrule
1 & 329\,ms & 3\,579\,384 & 22 & 3\,625\,773 & 7\,205\,179 & 85.6\% \\
2 & 291\,ms & 3\,899\,790 & 22 & 3\,988\,143 & 7\,887\,955 & 84.2\% \\
3 & 303\,ms & 4\,726\,689 & 22 & 4\,799\,899 & 9\,526\,610 & 80.9\% \\
\bottomrule
\end{tabular}
\end{center}

\begin{figure}[H]
  \centering
  \fig{task1b.png}{\linewidth}
  \caption{Task 1b: execution time and peak memory of each version, with 16 threads.}
\end{figure}

\subsection{Task 1c: producer and consumers}

Task~1c divides the work in a different way. A single thread, the reader, reads \texttt{numbers.txt}
line by line with a \texttt{StreamReader} and collects the lines into batches of 10\,000. Each full
batch is added to a thread-safe \texttt{BlockingCollection} whose capacity is limited to 100 batches.
When the queue is full, the call that adds a batch blocks the reader until a worker has taken a
batch out. The workers repeatedly take the next batch through \texttt{GetConsumingEnumerable}, which
blocks a worker while the queue is empty, and count the numbers of the batch in their own local
counters. When the reader reaches the end of the file, it calls \texttt{CompleteAdding} to close the
queue; the workers then finish as soon as the remaining batches have been taken, and their counts
are added together.

\begin{figure}[H]
  \centering
  \fig{task1c.png}{0.8\linewidth}
  \caption{Task 1c: execution time and speed-up against the number of workers.}
\end{figure}

\begin{center}
\small
\begin{tabular}{lrrrrrrr}
\toprule
Workers & 1 & 2 & 4 & 8 & 16 & 32 & 64 \\
\midrule
Time, ms (median)    & 2\,684 & \textbf{2\,630} & 3\,052 & 3\,430 & 4\,100 & 5\,453 & 4\,160 \\
Speed-up             & 1.00 & 1.02 & 0.88 & 0.78 & 0.65 & 0.49 & 0.65 \\
Peak working set, MB & 50 & 59 & 78 & 120 & 161 & 202 & 291 \\
\bottomrule
\end{tabular}
\end{center}

\begin{figure}[H]
  \centering
  \fig{task1a_vs_1c.png}{0.72\linewidth}
  \caption{Execution time of Task 1a and Task 1c against the number of threads.}
\end{figure}

\subsection{Answers to the questions of Task 1}

\question{1. At which number of threads does the speed-up stop growing? Relate it to your physical
cores and logical processors.}

In Task~1a, the speed-up keeps growing up to 16 threads, where it reaches 6.32\timesx, and then it
stops. With 32 threads it is 6.55\timesx{} and with 64 threads 6.01\timesx. These last differences
are smaller than the normal variation between runs (the runs with 64 threads alone took between 156
and 238\,ms), so the curve is simply flat after 16 threads. This is exactly the number of logical
processors of the laptop: with 16 threads, every hardware thread of the processor has work to do, and
there is no free hardware left.

The shape of the curve follows the structure of the processor. Up to 8 threads, every thread can have
a physical core to itself, and the speed-up grows to 4.58\timesx. The growth is not perfect, because
the efficiency falls from 89\% with 2 threads to 80\% with 4 and 57\% with 8. Between 8 and 16
threads, the threads have to share the cores through SMT, and the time falls only from 207 to
150\,ms, a further 1.38\timesx{} instead of 2\timesx. This is the expected behaviour: the two threads of a
core share its execution units, so the second one only fills the gaps in which the first one waits
(a parser has many such gaps, for example mispredicted branches and memory accesses). A gain of 38
per cent is somewhat higher than the 15 to 30 per cent mentioned in lecture~1, but of the same order.

The speed-up is therefore far below 16\timesx, and the likely reasons are the following. Only 8 of the 16
hardware threads are full cores. When all cores are busy, the processor also runs them at a lower
clock frequency than a single busy core, because of its power and temperature limits, so every
thread becomes slower as more threads are added. All threads copy their data out of the same file
cache, so they share the memory bandwidth and the last-level cache. Finally, the run takes only about
150\,ms at its best, so fixed costs such as creating the threads, opening the file once per thread and
the sequential steps at the beginning and at the end are no longer negligible. This last effect is
described by Amdahl's law.

\question{2. Why do 32 and 64 threads not help? Are they slower, and why?}

Once there are more threads than the 16 logical processors, there is no additional hardware on which
the extra threads could run. The operating system can only take turns, letting each thread run for a
short time slice before switching to another one. The total amount of CPU work stays exactly the
same, and on top of it come the costs of switching between threads and of refilling the processor
caches after each switch, so more threads than logical processors can never make a CPU-bound task
faster.

In Task~1a, 32 and 64 threads are nevertheless hardly different from 16 (145 and 158\,ms against
150\,ms). The threads of Task~1a share no data, and even with 64 threads every part of the file is
still about 5.5\,MB large, so the time lost to switching is small compared with the useful work. The
64-thread configuration is the slowest of the three, but the difference is within the variation
between runs (156 to 238\,ms for 64 threads alone), so it should not be over-interpreted. The extra
threads still cost memory, as the next question shows, and they bring no gain at all. In Task~1c,
where the threads compete for a shared queue, more threads are clearly slower, which is discussed in
question~5.

\question{3. Does the peak memory depend on the number of threads? Why?}

The peak memory depends on the number of threads whenever each thread owns some data of its own,
and it does not depend on it when the data is shared. Task~1a shows the first case very clearly.
Its peak private memory grows from 8\,MB with one thread to 73\,MB with 64 threads, which is about
1\,MB for every additional thread. That matches the 1\,MB read buffer that each thread allocates,
plus the small part of the thread's own stack that is actually used.

To confirm that the read buffer is responsible, Task~1a with 64 threads was run once more with a
buffer of 64\,KB per thread instead of 1\,MB, together with the normal 1\,MB configuration in the same
session (three runs of each, medians):

\begin{lstlisting}
.\bin\Release\net9.0\Labs.exe lab1 1a --threads 64                  # 1 MB buffer per thread (default)
.\bin\Release\net9.0\Labs.exe lab1 1a --threads 64 --buffer-kb 64   # 64 KB buffer per thread
\end{lstlisting}

\begin{center}
\small
\begin{tabular}{lrrrc}
\toprule
Task 1a, 64 threads & Time (median of 3) & Peak working set & Peak private & Check \\
\midrule
1\,MB buffer per thread  & 99\,ms & 91\,MB & 74\,MB & OK \\
64\,KB buffer per thread & 88\,ms & 31\,MB & 15\,MB & OK \\
\bottomrule
\end{tabular}
\end{center}

With the smaller buffers, the peak private memory fell from 74\,MB to 15\,MB and the peak working
set from 91\,MB to 31\,MB, which is only about 6\,MB above the 25\,MB that the .NET runtime occupies
before any work starts. The saving of about 60\,MB is exactly what the buffers predict:
$64 \times (1024 - 64)\,\text{KB} = 60\,\text{MB}$. The run did not become slower (88\,ms against
99\,ms, a difference within the variation between runs), because a 64\,KB read is still large enough
that the fixed cost of each read call is negligible compared with parsing 64\,KB of numbers. The
private memory per thread is now about 100\,KB, the 64\,KB buffer plus the used part of the stack.
These two configurations were measured later than the main table of Task~1a, with the laptop in a
faster state, so their times should only be compared with each other.

Task~1c also uses more memory with more workers, growing from about 50\,MB to 291\,MB in the largest
configuration. Here the memory is mostly made of the strings that the reader creates. Each worker
holds a batch of 10\,000 strings while it counts it, so with more workers more batches are in use at
the same time. When the garbage collector runs, these batches are still alive, so they survive the
collection and are moved to an older generation of the heap, which is cleaned up only rarely.

Task~1b shows the opposite case. Its memory is the same in every version, about 218\,MB, because it
is dominated by the shared array of 50 million integers (200\,MB), while the counters add only a few
bytes per thread. Task~2 again grows with the number of threads, because each thread holds a whole
44\,MB file in memory, as described in Section~\ref{sec:task2}. In summary, memory grows with the
number of threads only for data that every thread owns privately, such as buffers and batches, while
shared, read-only data costs the same regardless of how many threads read it.

\question{4. Rank the three versions of Task 1b by time. Which of them are correct? Explain each
result.}

Ordered from the fastest to the slowest, the versions are local (77\,ms), unsynchronised (302\,ms),
interlocked (1\,057\,ms) and locked (10\,261\,ms). The local, interlocked and locked versions produce
the correct totals. The unsynchronised version does not.

The unsynchronised version is wrong because an increment such as \texttt{shared.Positive++} is not a
single step. The processor has to read the current value from memory, add one to it, and write the
result back. If two threads read the same old value at nearly the same time, both add one and both
write back the same new value, so one of the two increments is simply lost. With 16 threads
constantly incrementing the same counters, this happened for between about 81 and 86 per cent of all
increments, and because it depends on the exact timing of the threads, the result was different in
every run (between 7.2 and 9.5 million counted instead of 50 million). The count of zeros was correct
every time, but only by luck: there are only 22 zeros among 50 million numbers, so two threads almost
never increment that counter at the same moment. Interestingly, the unsynchronised version is also
about 4 times slower than the local one, even though it does no synchronisation at all. The three
shared counters lie in the same 64-byte cache line, and each write by one core invalidates that line
in the caches of the other cores, so the line is constantly passed back and forth between the cores.

The locked version is correct, but it is the slowest by far, about 130 times slower than the local
one. Each of the 50 million increments has to acquire and release the lock, so only one thread can
make progress at any moment, while the other 15 wait. The work inside the lock is therefore performed
one thread at a time again, exactly as lecture~2 described the lock as a queue, and on top of that
every hand-over of the lock between threads costs time, about 205\,ns of wall-clock time per number.

The interlocked version is also correct and is about 14 times slower than the local one.
\texttt{Interlocked\allowbreak.Increment} makes the read, the addition and the write a single atomic hardware
operation, so no update can be lost and no thread is ever blocked. Each atomic increment still needs
exclusive ownership of the cache line that holds the counter, however, so the cores still have to
take turns, only at a much lower cost than with a lock (about 21\,ns per number).

The local version is both correct and the fastest. Each thread counts in its own local variables,
which the compiler keeps in processor registers, so there is no shared state while counting, and
therefore neither a race nor any waiting. Each thread writes its result only once, at the very end,
and the main thread adds the 16 results after \texttt{Join} has guaranteed that all threads are
finished. This is the first step of the ladder of fixes from lecture~2, ``don't share'', and the
measurement confirms its main message: it is correct and at the same time even faster than the
broken version.

\question{5. Compare Task 1c with Task 1a. Which one is faster, and why? At which number of workers
does Task 1c stop getting faster?}

Task~1a is faster than Task~1c for every number of threads. With one thread, Task~1a needs 949\,ms
and Task~1c 2\,684\,ms, and at their best configurations Task~1a needs 145\,ms while Task~1c needs
2\,630\,ms, which makes Task~1a about 18 times faster.

The main reason is that Task~1c contains a stage that cannot be parallelised. A single reader thread
has to read and split every one of the 50 million lines, so no matter how many workers parse the
numbers, the whole run can never be shorter than the time the reader alone needs, which is about
2.6\,s. This is an example of Amdahl's law: the sequential part of a program limits its total
speed-up. In Task~1a, by contrast, both the reading and the parsing are divided among all threads.
In addition, the reader of Task~1c does more work for each line than Task~1a does. It converts the
bytes into characters and creates a separate string object for each line, which means 50 million
objects that the garbage collector afterwards has to clean up, and the workers then parse these
strings with \texttt{int.Parse}. Task~1a reads the digits directly from the byte buffer and creates
no objects at all. Passing the work through the queue also has a cost, although it is small with
batches of 10\,000 lines, because there are only about 5\,000 queue operations in total.

Task~1c stops getting faster immediately. With one worker it takes 2\,684\,ms and with two workers
2\,630\,ms, an improvement of only 2 per cent. A single worker already keeps up with the reader almost
completely, so any additional workers mostly wait on an empty queue. From four workers on, the task
becomes slower, reaching 3\,052\,ms with 4 workers and more than 4\,000\,ms with 16 or more. More
workers compete for the internal lock of the queue, more threads have to be put to sleep and woken up
again, and the garbage collector has more threads to stop every time it runs. The values for 32 and
64 workers (5\,453 and 4\,160\,ms) are not in the expected order, which is measurement noise: the
runs with 32 workers alone ranged from about 4.8 to 5.8\,s.

\question{6. What would happen with batches of a single line? And with a queue that has no size
limit?}

Both variants were measured with 1 and with 8 workers, together with the normal configuration of the
assignment in the same session, so that the three rows can be compared directly (three runs of each,
medians; every run reported \OK). The commands were:

\begin{lstlisting}
.\bin\Release\net9.0\Labs.exe lab1 1c --workers 8                  # the assignment: batches of 10000, at most 100
.\bin\Release\net9.0\Labs.exe lab1 1c --workers 8 --batch 1        # batches of a single line
.\bin\Release\net9.0\Labs.exe lab1 1c --workers 8 --capacity 0     # a queue with no size limit
\end{lstlisting}

and the same three commands with \texttt{-{}-workers 1}.

\begin{center}
\small
\begin{tabularx}{\linewidth}{@{}Xrrr@{}}
\toprule
Variant & 1 worker & 8 workers & Peak working set (1 / 8 workers) \\
\midrule
batches of 10\,000 lines, at most 100 batches (the assignment) & 1\,359\,ms & 1\,724\,ms & 51 / 120\,MB \\
\textbf{batches of 1 line}, at most 100 batches & 14\,841\,ms & 43\,959\,ms & 34 / 34\,MB \\
batches of 10\,000 lines, \textbf{no size limit} & 1\,353\,ms & 1\,764\,ms & 50 / 120\,MB \\
\bottomrule
\end{tabularx}
\end{center}

These runs were made later than the main table of Task~1c, with the laptop in a faster state (the
normal configuration took 1.36 and 1.72\,s instead of 2.68 and 3.43\,s), so they should be compared
with each other and not with that table.

With batches of a single line, every number becomes a separate queue operation, so the reader and
the workers have to take the queue's internal lock and signal one another 50 million times instead of
about 5\,000 times. This coordination costs far more than counting a single number, and the
measurement shows it clearly. With one worker, the run took 14.8\,s, about 11 times longer than with
batches of 10\,000 lines. With eight workers it took 44.0\,s (the three runs ranged from 38.2 to
47.6\,s), about 26 times longer than the normal configuration and three times longer than with a
single worker. More workers therefore make it worse, because all of them fight over the same lock for
every single number and are repeatedly put to sleep and woken up again whenever the queue runs empty.
With eight workers, each number cost about 0.9\,\textmu s of wall-clock time, far more than the parsing
and counting itself. The memory, on the other hand, was the lowest of all Task~1c runs: a peak working
set of 34\,MB with any number of workers, only about 9\,MB above what the runtime occupies on its own,
because the queue can never hold more than 100 lines and no worker holds a batch of 10\,000 strings.
The experiment shows the purpose of batching: each queue operation has a fixed cost, and a batch must
contain enough work to make that cost negligible.

A queue without a size limit removes the mechanism that slows the reader down when the workers
cannot keep up. The reader never has to wait, and this missing feedback is called a lack of
backpressure. On this laptop the change made no measurable difference: 1.35\,s against 1.36\,s with
one worker and 1.76\,s against 1.72\,s with eight, which is within the variation between runs, and the
peak memory was the same (50 and 120\,MB). This confirms that the reader is the slow side: the
workers empty the queue faster than the reader fills it, so the queue never comes close to 100
batches, and the limit never has to block the reader. If the workers were slower, however, for
example because the work per number were heavier, the queue would keep growing until it held the
entire file as strings. That would be roughly 50 million strings of 40 to 48 bytes each, or about 2
to 2.4\,GB, compared with the roughly 48\,MB that 100 batches can hold. The size limit of the queue is
therefore what keeps the memory bounded: it forces the reader to adapt to the speed of the workers,
and not the other way round.

% Section 4
\section{Task 2: Many files on slow storage}\label{sec:task2}

In Task~2, a whole file is the unit of work. A thread reads a file in a single operation with
\texttt{File.ReadAllBytes}, which places the entire file in one byte buffer, and then parses the
buffer with the same parser that Task~1a uses. The 8 files are dealt to the threads in turn (file
$i$ goes to thread $i \bmod T$), so that each thread receives the same number of files, $8/T$.
Besides the total time, each thread also measures how long its reading and its parsing take, which
makes it possible to see where the time is spent. The files are read only from the USB flash drive,
whose volume was dismounted before every run so that nothing came from the file cache. An earlier
attempt without this step read all 8 files in only 468\,ms, about 800\,MB/s, which is far too fast for
a flash drive: that data came from the cache, and the run is not used in this report. The runs below
were started with the following commands, each one preceded by a dismount (three runs of each,
medians):

\begin{lstlisting}
.\lab1\scripts\dismount-volume.ps1 -Drive D
.\bin\Release\net9.0\Labs.exe lab1 2 --dir D:\lab1 --threads 1 --label usb-cold --csv lab1\results\task2.csv
# the same with --threads 2 and --threads 4
\end{lstlisting}

\begin{figure}[H]
  \centering
  \fig{task2.png}{\linewidth}
  \caption{Task 2: execution time, speed-up and peak working set against the number of threads,
           with the files on the USB drive and an empty cache.}
\end{figure}

\begin{center}
\small
\begin{tabular}{lrrr}
\toprule
Threads & 1 & 2 & 4 \\
\midrule
Time (median of 3)            & 11\,391\,ms & 10\,960\,ms & 10\,908\,ms \\
Speed-up                     & 1.00 & 1.04 & 1.04 \\
Reading, summed over threads & 10\,959\,ms & 20\,118\,ms & 35\,137\,ms \\
Parsing, summed over threads & 428\,ms & 425\,ms & 423\,ms \\
Peak working set             & 246\,MB & 246\,MB & 379\,MB \\
Peak private                 & 228\,MB & 228\,MB & 361\,MB \\
Check                        & OK & OK & OK \\
\bottomrule
\end{tabular}
\end{center}

\begin{figure}[H]
  \centering
  \fig{task2_breakdown.png}{0.85\linewidth}
  \caption{Where the time goes with one thread: Task 2 reading from the USB drive, compared with
           Task 1a reading the cached file.}
\end{figure}

\subsection{Answers to the questions of Task 2}

\question{1. Compare the time on one thread with Task 1a on one thread. On the USB drive, is the work
CPU-bound or I/O-bound?}

Task~1a processed exactly the same numbers with one thread in 949\,ms, reading from the file cache.
Task~2 with one thread took 11\,391\,ms, so it took about 12 times longer. The time measured inside
each thread shows where the difference comes from: reading the 8 files took 10\,959\,ms, which is
96\% of the total and corresponds to about 32\,MB/s (352\,MB in 10.96\,s), while parsing the numbers
took only 428\,ms. Parsing is the same CPU work as in Task~1a, and it is even shorter here because the
whole file is already in one buffer and no time is spent copying it out of the cache. For 96\% of the
run the processor is idle and waits for the drive, so on the USB drive the work is clearly I/O-bound.
A faster processor could save at most the 4\% parsing share, whereas a faster drive would shorten the
run almost in proportion to its speed.

\question{2. On the USB drive, do more threads help? Why not, when lecture 1 showed that I/O-bound
work benefits from concurrency?}

On the flash drive, more threads hardly help: 2 threads gave a speed-up of 1.04 (10\,960\,ms) and 4
threads also 1.04 (10\,908\,ms). A single flash drive connected through a single USB link delivers a
fixed amount of data per second, however many requests it receives at the same time. Parallel reads
simply share this bandwidth between them, which the measurements show directly: the reading time
summed over all threads grows from 11.0\,s with one thread to 20.1\,s with two and 35.1\,s with four,
because every thread now spends most of its time waiting while the drive serves the other threads,
whereas the total time stays the same. The only part of the work that can overlap is the parsing,
whose share of the total time is $p$ (so the reading share is $1-p$). Amdahl's law limits the
speed-up with $T$ threads to
\[
  S(T) = \frac{1}{(1-p) + p/T}, \qquad S(\infty) = \frac{1}{1-p}.
\]
With the measured $p = 428 / 11\,391 \approx 0.038$, this gives $S(2) \approx 1.02$,
$S(4) \approx 1.03$ and an absolute limit of $S(\infty) \approx 1.04$, however many threads are used.
The measurement reaches this limit already with two threads, because while one thread parses a file,
another one is already reading the next one, so the parsing is completely hidden behind the reading.
The run with two threads (10.96\,s) lasts exactly as long as the pure reading with one thread
(10.96\,s): the drive was busy all the time, and nothing more can be gained.

Lecture~1 is nevertheless right that concurrency helps I/O-bound work, but only for a particular kind
of waiting. When a program waits for a server on the network, most of the waiting is latency: the
request travels to the server, the remote machine does its work, and the answer travels back, while
the local network connection is almost idle the whole time. Ten such requests can wait at the same
time, so a concurrent program overlaps their waiting periods, and the total time drops almost
tenfold. Concurrency helps when the waiting is spent on someone else's capacity, which is called
latency-bound I/O. It does not help when the waiting is caused by the limited bandwidth of the device
itself, which is called throughput-bound I/O, as with a single USB stick: a resource that is already
fully busy cannot deliver data any faster just because more threads are asking for it.

\textbf{Memory.} The peak memory of Task~2 also grows with the number of threads (a peak working set
of 246\,MB with one and two threads and 379\,MB with four), because every thread holds a complete
44\,MB file in a single \texttt{byte[]}. Even with one thread, the peak private memory of 228\,MB
corresponds to about five files, and with four threads 361\,MB corresponds to all eight. Arrays of this size are placed on the large object heap of
.NET, which is cleaned up only during the rare full (generation 2) garbage collections, so several
buffers of files that have already been processed remain in memory until the next such collection
runs.

% Section 5
\section{Conclusion}

The measurements show that on this laptop Task~1 is CPU-bound in all three of its forms, while Task~2
is I/O-bound. The table summarises the evidence for each part.

\begin{center}
\small
\begin{tabularx}{\linewidth}{@{}p{3.3cm}p{3.1cm}X@{}}
\toprule
Work & Limited by & Evidence from the measurements \\
\midrule
Task 1a (parsing a cached file) & the CPU &
  The time falls as threads are added, up to 6.32\timesx{} at 16 threads, and stops improving at the
  16 logical processors (8 physical cores with SMT). \\
Task 1b (counting an in-memory array) & the CPU and the synchronisation &
  There is no I/O at all. The time depends only on how the threads share the counters: 77\,ms with
  local counters against 10\,261\,ms with a lock. \\
Task 1c (one reader, $N$ workers) & the CPU of a single thread &
  The file is in the RAM cache (Task 1a reads it in 145\,ms at best), yet Task 1c needs about
  2.6\,s with any number of workers, because the single reader, which creates 50 million strings,
  sets the pace. \\
Task 2 (8 files on a USB flash drive) & the bandwidth of the drive &
  96\% of the time is spent reading, at about 32\,MB/s. Two and four threads give a speed-up of only
  1.04, while the reading time summed over the threads grows from 11.0 to 35.1\,s. \\
\bottomrule
\end{tabularx}
\end{center}

The two kinds of work can be told apart from the measurements alone. CPU-bound work becomes faster
when threads are added, until the processor runs out of hardware threads, which on this machine
happens at 16 threads, and during the whole run the processor is busy. I/O-bound work leaves the
processor mostly idle, and its duration is set by the device it waits for. Additional threads help
I/O-bound work only if the device can serve several requests in parallel, which a single USB flash
drive cannot do.

The laboratory work also illustrates a lesson about correctness. Shared counters without
synchronisation are fast but wrong, and protecting them with a lock makes them correct but 130 times
slower. The best design is not to share the counters at all, but to let every thread count on its own
and to add the results after \texttt{Join}, which is both correct and the fastest of all versions.
The correctness check proved its value here: it detected the race condition in the unsynchronised
version of Task~1b, which still produced plausible-looking numbers that were simply wrong.

% Appendices
\appendix

\section{Building and running the program}\label{app:run}

The program requires the .NET 9 SDK, Windows PowerShell 5.1 or newer for the measuring scripts, and
about 1\,GB of free disk space for the generated data. The commands below are run in PowerShell; the
options of the program are written exactly as the program expects them.

\begin{lstlisting}
# from the repository root (the folder with Labs.csproj)
dotnet build Labs.csproj -c Release
.\bin\Release\net9.0\Labs.exe lab1 info          # hardware report
.\bin\Release\net9.0\Labs.exe lab1 gen           # numbers.txt and the reference counts
.\bin\Release\net9.0\Labs.exe lab1 split         # 8 parts for Task 2

# single runs
.\bin\Release\net9.0\Labs.exe lab1 1a --threads 8
.\bin\Release\net9.0\Labs.exe lab1 1b --mode unsync      # unsync, lock, local or interlocked
.\bin\Release\net9.0\Labs.exe lab1 1c --workers 4        # batches of 10000, at most 100 in the queue
.\bin\Release\net9.0\Labs.exe lab1 2 --dir E:\lab1 --threads 2

# additional experiments (Task 1, questions 3 and 6)
.\bin\Release\net9.0\Labs.exe lab1 1a --threads 64 --buffer-kb 64   # 64 KB buffer per thread instead of 1 MB
.\bin\Release\net9.0\Labs.exe lab1 1c --workers 8 --batch 1         # batches of a single line
.\bin\Release\net9.0\Labs.exe lab1 1c --workers 8 --capacity 0      # a queue with no size limit
.\bin\Release\net9.0\Labs.exe lab1 1c --workers 1 --batch 1         # the same two variants with one worker
.\bin\Release\net9.0\Labs.exe lab1 1c --workers 1 --capacity 0

# all measurements (from the lab1\scripts folder)
Set-ExecutionPolicy -Scope Process Bypass
.\run-task1.ps1                              # Task 1a, 1b and 1c, 3 repetitions each
.\run-task2.ps1 -UsbDir E:\lab1 -CopyFirst   # asks for an eject and replug before each run
.\dismount-volume.ps1 -Drive D               # or: empty the drive's cache without unplugging it
.\summary.ps1                                # medians and speed-ups in results\summary.csv
python ..\report\make_charts.py              # charts in report\figures
\end{lstlisting}

\section{Raw data}\label{app:data}

All measured runs are stored as CSV files in the \texttt{lab1/results} folder of the source code.
The files \texttt{task1a.csv}, \texttt{task1b.csv}, \texttt{task1c.csv} and \texttt{task2.csv}
contain every individual run with its time, memory, counts and the result of the correctness check.
The file \texttt{summary.csv} contains, for every configuration, the median, minimum and maximum
time, the speed-up and the median memory. The tables and charts of Task~1 use the three repetitions
of 4 October 2026; \texttt{make\_charts.py} selects them by date, because the CSV files also contain
a later set of repetitions. In \texttt{task2.csv}, the runs used in this report have the variant
\texttt{usb-cold}; the single earlier run with the variant \texttt{usb} came from the file cache and
is not used. The console
output of the additional experiments of questions 3 and 6, which were started by hand and are
therefore not in these CSV files, is stored in \texttt{extra\_runs\_2026-10-13.txt}.

\end{document}
